\clearpage
\section{Rank-four normalization and induction in stages}\label{app:rank4_stages}

This is a corollary of the enlarged-carrier construction in
Appendix~\ref{app:rank4}.

Use that note's actual groups and notation:

\begin{Verbatim}[fontsize=\small,breaklines=true,breakanywhere=true]
n=2a,  r=2m+1,  a>r,
H subset J=J1(beta) normal in P=J U_B^m,
[P:J]=2,  N=K^times J,
pi=c-Ind_N^G(Lambda),  Lambda|J=eta.
\end{Verbatim}

The field K normalizes P. Consequently the following are actual groups
and actual smooth representations:

\begin{Verbatim}[fontsize=\small,breaklines=true,breakanywhere=true]
M=K^times P,
Sigma=Ind_N^M(Lambda),
dim Sigma=4,
Sigma|P=Ind_J^P(eta),
pi is isomorphic to c-Ind_M^G(Sigma).
\end{Verbatim}

The representation Sigma is irreducible, and M is compact modulo the
center of G. Thus the enlarged carrier has an exact arithmetic inducing
datum for the same candidate, not merely an occurrence in its
restriction.

\hypertarget{rank4_stages-normalization-by-the-actual-parameter-field}{%
\subsection{1. Normalization by the actual parameter
field}\label{rank4_stages-normalization-by-the-actual-parameter-field}}

The units O\_K\^{}times equal U\_K\^{}1 because the residue field is
F\_2. They belong to J and hence normalize P. It remains to check a
uniformizer k of K. This element has A-grade one and normalizes A.

The actual-field approximation estimate proved in Section~2 of the
carrier note gives

\begin{Verbatim}[fontsize=\small,breaklines=true,breakanywhere=true]
k=b+e,  b in Q^1,  e in P_A^(n-r+1).
\end{Verbatim}

The leading grade of b is the nonzero centralizing line of the primitive
derived parameter. Over F\_2 it is exactly the leading grade of R, the
period-two E-chain shift. Therefore

\begin{Verbatim}[fontsize=\small,breaklines=true,breakanywhere=true]
b=R v,  v in 1+Q^1.
\end{Verbatim}

In particular b is invertible and normalizes both the complete E-chain
and the complete interleaved F-chain. We can write

\begin{Verbatim}[fontsize=\small,breaklines=true,breakanywhere=true]
k=b(1+epsilon),  epsilon in P_A^(n-r).
\end{Verbatim}

For u in U\_B\^{}m, the usual filtration commutator bound gives

\begin{Verbatim}[fontsize=\small,breaklines=true,breakanywhere=true]
(1+epsilon)u(1+epsilon)^(-1)u^(-1)
   in U_A^(n-r+m).
\end{Verbatim}

This group is contained in U\_A\^{}a subset J, because

\begin{Verbatim}[fontsize=\small,breaklines=true,breakanywhere=true]
n-r+m=2a-m-1>a.
\end{Verbatim}

Conjugation by b preserves both relevant filtrations. Since b normalizes
U\_B\^{}m, it follows that

\begin{Verbatim}[fontsize=\small,breaklines=true,breakanywhere=true]
k U_B^m k^(-1) subset U_B^m J=P.
\end{Verbatim}

Also kJk\^{}(-1)=J, as K\^{}times normalizes its actual simple-character
groups. Thus kPk\^{}(-1) is contained in P and contains J. Its index
over J is two, equal to {[}P:J{]}, so it equals P. This proves
normalization by k, its inverse, and hence all K\^{}times.

No equality or inclusion between the fields E and K was used. The
element b belongs to the actual outer centralizer B and only
approximates k to the stated precision.

\hypertarget{rank4_stages-the-two-inducing-groups-have-index-two}{%
\subsection{2. The two inducing groups have index
two}\label{rank4_stages-the-two-inducing-groups-have-index-two}}

Since K\^{}times normalizes P, the product M=K\^{}times P is a group.
Both P/J and the conjugation action of K\^{}times on it are defined; P/J
has order two, whose automorphism group is trivial. Therefore

\begin{Verbatim}[fontsize=\small,breaklines=true,breakanywhere=true]
[K^times,P] subset J.
\end{Verbatim}

Together with J normal in P and K\^{}times normalizing J, this proves
N=K\^{}times J normal in M. The intersection K\^{}times cap P is
U\_K\^{}1, which is contained in J, since P subset U\_A\^{}1.
Consequently

\begin{Verbatim}[fontsize=\small,breaklines=true,breakanywhere=true]
[M:N]=[P:J]=2.
\end{Verbatim}

The quotient \(K^\times/F^\times\) is compact because K/F is a
finite totally ramified extension, and P is compact. Hence M is compact
modulo F\^{}times. It is open in G because it contains the open subgroup
J.

\hypertarget{rank4_stages-the-enlarged-inducing-representation-is-irreducible}{%
\subsection{3. The enlarged inducing representation is
irreducible}\label{rank4_stages-the-enlarged-inducing-representation-is-irreducible}}

Let u in P represent its nontrivial coset modulo J. The carrier note
proved that theta\^{}u differs from theta, and therefore eta\^{}u
differs from eta. It follows that Lambda\^{}u differs from Lambda, since
their restrictions to J already differ. Applying finite-index Clifford
theory to N normal in M gives

\begin{Verbatim}[fontsize=\small,breaklines=true,breakanywhere=true]
Sigma=Ind_N^M(Lambda) irreducible,
dim Sigma=2 dim Lambda=4.
\end{Verbatim}

There is one N--P double coset in M, and N cap P=J. Thus restriction to
P gives the actual four-dimensional Heisenberg carrier:

\begin{Verbatim}[fontsize=\small,breaklines=true,breakanywhere=true]
Sigma|P is isomorphic to Ind_J^P(eta).
\end{Verbatim}

Its central character on F\^{}times is unchanged from Lambda, because
F\^{}times is central and lies in N.

\hypertarget{rank4_stages-exact-recovery-by-induction-in-stages}{%
\subsection{4. Exact recovery by induction in
stages}\label{rank4_stages-exact-recovery-by-induction-in-stages}}

Induction in stages for these open subgroups gives

\begin{Verbatim}[fontsize=\small,breaklines=true,breakanywhere=true]
c-Ind_M^G(Sigma)
   is isomorphic to c-Ind_M^G(Ind_N^M Lambda)
   is isomorphic to c-Ind_N^G(Lambda)
   =pi.
\end{Verbatim}

The second induction is finite-index, so ordinary and compact induction
from N to M agree. This is an exact equality of the resulting
isomorphism classes and applies to the same fixed Lambda, including its
noncompact scalar choice. It introduces no new candidate and no H5
premise.

\hypertarget{rank4_stages-scope}{%
\subsection{Scope}\label{rank4_stages-scope}}

This establishes the arithmetic inducing datum and exact recovery for
the enlarged carrier. It leaves unchanged the distinction between the
original two-dimensional BK quotient J/H and the new four-dimensional
actual quotient P/H0. It also leaves unchanged the source-phase and
literal extraspecial-center limitations recorded in the carrier note.

Dependency pin: \path{RANK4_ENLARGED_CARRIER167.md} SHA256
\path{1517e26138776fe0f32d74a174d309ec5dd6a0f572a979473d8af3d20b7a4155}.
